\section{Acceptance criterion ablation (exploratory)}
\label{sec:ablation}

\subsection{Motivation}

The R2 results of Section~\ref{sec:mismatch} mix two effects:

\begin{enumerate}
  \item the initial plan going stale as the minibatch changes;
  \item steps being rejected because monotone decrease is required on a noisy
        objective.
\end{enumerate}

The default acceptance rule requires monotone decrease of the current minibatch
objective. A step that improves the true objective can be rejected because of
sample noise.

\subsection{Design}

\begin{quote}
We replaced minibatch-local monotonic acceptance with a fixed-evaluation-objective
criterion.
\end{quote}

Gradients and HVPs are still computed on the minibatch. Only the acceptance test
changes. A forward pass on the fixed evaluation objective costs
$n_{\text{samples}} / \text{batch size}$ times more, and that factor is included in
the GE accounting.

\subsection{Results}

\begin{table}[!ht]
  \centering
  % 이 파일은 scripts/make_tables.py 가 생성한다. **손으로 수정하지 않는다.**
% 표: acceptance criterion ablation (exploratory, n=3 per regime)
% raw: headroom_micro-neural_step_size_fixed_b8_0bec1125.jsonl
% raw: headroom_micro-neural_step_size_fixed_b8_9f3194be.jsonl
% raw 의 SHA-256 은 paper/evidence_map.md 에 있다.
% 공개 저장소에는 raw 가 없다. results/public/*.csv 에서 --public-dir 로 만들면
% **고정된 환경에서** 이 파일과 바이트 단위로 같아진다.
% 판본이 다른 환경에서의 비트 수준 일치는 주장하지 않는다.
\begin{tabular}{lllrr}
\toprule
Gate & Regime & & Control & Fixed evaluation \\
\midrule
A2 & R1 full-batch & & $+15.176$ (3/3) & $+14.065$ (3/3) \\
 & R2 batch 128 & & $-0.900$ (1/3) & $-0.598$ (0/3) \\
 & R2 batch 64 & & $-0.277$ (1/3) & $-0.093$ (0/3) \\
\midrule
C2 & R1 full-batch & & $+0.547$ (2/3) & $-0.080$ (1/3) \\
 & R2 batch 128 & & $-1.111$ (1/3) & $-0.118$ (1/3) \\
 & R2 batch 64 & & $-0.716$ (1/3) & $+0.104$ (2/3) \\
\midrule
C3 & R1 full-batch & & $-0.095$ (1/3) & $-0.949$ (0/3) \\
 & R2 batch 128 & & $+2.941$ (3/3) & $+1.110$ (3/3) \\
 & R2 batch 64 & & $+1.666$ (3/3) & $+0.973$ (3/3) \\
\bottomrule
\end{tabular}

  \caption{Paired differences under the two acceptance criteria, $n = 3$ per
  regime. Exploratory.}
  \label{tab:acceptance-delta}
\end{table}

\begin{figure}[!ht]
  \centering
  \includegraphics[width=\textwidth]{figure4_acceptance_ablation.png}
  \caption{Paired \gate{A2}, \gate{C2} and \gate{C3} effects under minibatch-local
  and fixed-evaluation acceptance, $n = 3$ per regime. The alternative criterion
  reduced, but did not eliminate, the apparent advantage of replanning over a
  committed stale plan, and \ctrl{shrinking} did not establish an advantage over
  \ctrl{onestep}. Exploratory.}
  \label{fig:acceptance}
\end{figure}

\begin{quote}
This alternative criterion reduced, but did not eliminate, the apparent advantage
of replanning over a committed stale plan. The estimated $\gate{C3}$ magnitude
decreased substantially under the alternative acceptance criterion, indicating that
the original magnitude was partly acceptance-dependent.
\end{quote}

In R2, \ctrl{shrinking} still does not beat \ctrl{best\_static} (zero of three
pairs positive) and only draws level with \ctrl{onestep}.

\subsection{The alternative criterion was stricter, not looser}
\label{sec:stricter}

\begin{table}[!ht]
  \centering
  % 이 파일은 scripts/make_tables.py 가 생성한다. **손으로 수정하지 않는다.**
% 표: rejection rate under two acceptance criteria
% raw: headroom_micro-neural_step_size_fixed_b8_0bec1125.jsonl
% raw: headroom_micro-neural_step_size_fixed_b8_9f3194be.jsonl
% raw 의 SHA-256 은 paper/evidence_map.md 에 있다.
% 공개 저장소에는 raw 가 없다. results/public/*.csv 에서 --public-dir 로 만들면
% 이 파일과 바이트 단위로 같아진다 (scripts/verify_public_results.py).
\begin{tabular}{llrr}
\toprule
Controller & Regime & Control & Fixed evaluation \\
\midrule
\texttt{committed} & R1 full-batch & 0.00 & 0.04 \\
 & R2 batch 128 & 0.79 & 0.89 \\
 & R2 batch 64 & 0.66 & 0.92 \\
\midrule
\texttt{shrinking} & R1 full-batch & 0.00 & 0.04 \\
 & R2 batch 128 & 0.04 & 0.42 \\
 & R2 batch 64 & 0.00 & 0.36 \\
\midrule
\texttt{onestep} & R1 full-batch & 0.04 & 0.04 \\
 & R2 batch 128 & 0.03 & 0.53 \\
 & R2 batch 64 & 0.00 & 0.57 \\
\bottomrule
\end{tabular}

  \caption{Median rejection rate under the two acceptance criteria. Rejection rates
  \emph{rose} under the alternative criterion.}
  \label{tab:acceptance-rejection}
\end{table}

Under the default rule, \texttt{loss\_before} and the candidate losses live on
\emph{the same minibatch}. Reducing the loss of the batch whose gradient was just
computed is easy. The alternative criterion requires a decrease in the full-data
loss and therefore blocks steps that overfit the current batch.

\subsection{Limitations of this ablation}
\label{sec:ablation-limits}

\textbf{This is not a single-factor change.} The acceptance criterion and the
number of effective steps change together. Excluding the evaluation cost from the
accounting would hide a real cost, so we included it, and the price is that the two
effects are entangled. We therefore \emph{do not use absolute comparisons between
the two columns to argue that one acceptance rule is better}.

The R1 column is not a meaningful comparison. In the full-batch regime the fixed
evaluation objective and the curvature objective take the same value, so all the
alternative criterion changes is charging one extra forward pass per step
(\ctrl{best\_static}: $5.2199 \rightarrow 5.1094$). That is a budget perturbation.
